\documentclass[12pt,letterpaper]{exam}
\usepackage[lmargin=1in,rmargin=1in,tmargin=1in,bmargin=1in]{geometry}
\usepackage{../style/exams}

% -------------------
% Course & Exam Information
% -------------------
\newcommand{\course}{MATH 142: Exam 2}
\newcommand{\term}{Spring ---\textsubscript{\textsubscript{2}} 2026}
\newcommand{\examdate}{03/19/2026}
\newcommand{\timelimit}{75 Minutes}

\setbool{hideans}{false} % Student: True; Instructor: False

\usepackage{cancel}

% -------------------
% Content
% -------------------
\begin{document}

\examtitle
\instructions{Write your name on the appropriate line on the exam cover sheet. This exam contains \numpages\ pages (including this cover page) and \numquestions\ questions. Check that you have every page of the exam. Answer the questions in the spaces provided on the question sheets. Be sure to answer every part of each question and show all your work and fully justify all your responses. If you run out of room for an answer, continue on the back of the page --- being sure to indicate the problem number.
} 
\scores
\bottomline
\newpage


% -------------------
% Questions
% -------------------
\begin{questions}

% Question 1
\newpage
\question[16] For each of the following series, determine whether the series converges or diverges. If the series converges, find the \textit{exact} value of the sum---if possible. \par\vspace{0.3cm}
	\begin{parts}
	% (a)
	\part $\ds\sum_{n=5}^\infty \cos \left( \dfrac{1}{n^4} \right)$ \par\vspace{0.5cm}
	\psol{\itshape We know\dots}
		\[
		\psol{\lim_{n \to \infty} \cos \left( \dfrac{1}{n^4} \right)= \cos(0)= 1 \neq 0}
		\]
	\psol{\itshape Therefore, the series $\ds\sum_{n=5}^\infty \cos \left( \dfrac{1}{n^4} \right)$ diverges by the Divergence Test.} \vfill
	
	% (b)
	\part $\ds\sum_{n=1}^\infty \big( \sqrt{n + 1} - \sqrt{n} \big)$ \par\vspace{0.3cm}
	\psol{\itshape This is a telescoping series. Observe\dots}
		\[
		\begin{aligned}
		\psol{\hspace{-2.75cm} \sum_{n=1}^\infty \big( \sqrt{n + 1} - \sqrt{n} \big)}&\psol{= \lim_{N \to \infty} \sum_{n=1}^N \big( \sqrt{n + 1} - \sqrt{n} \big)} \\[0.3cm]
		&\resizebox{.9\hsize}{!}{\psol{$\ds =\lim_{N \to \infty} \left[ \left( \sqrt{2} - \sqrt{1} \right) + \left( \sqrt{3} - \sqrt{2} \right) + \left( \sqrt{4} - \sqrt{3} \right) + \cdots + \left( \sqrt{N} - \sqrt{N - 1} \right) + \left( \sqrt{N + 1} - \sqrt{N} \right) \right]$}} \\[0.3cm]
		&\resizebox{.9\hsize}{!}{\psol{$\ds =\lim_{N \to \infty} \left[ \left( \cancel{\sqrt{2}} - \sqrt{1} \right) + \left( \cancel{\sqrt{3} - \sqrt{2}} \right) + \left( \cancel{\sqrt{4}} - \cancel{\sqrt{3}} \right) + \cdots + \left( \cancel{\sqrt{N}} - \cancel{\sqrt{N - 1}} \right) + \left( \sqrt{N + 1} - \cancel{\sqrt{N}} \right) \right]$}} \\[0.3cm]
		& \psol{= \lim_{N \to \infty} \left[ \sqrt{N + 1} - 1 \right]} \\[0.3cm]
		& \psol{= \infty}
		\end{aligned}
		\]
	\psol{\itshape Therefore, the series $\ds\sum_{n=1}^\infty \big( \sqrt{n + 1} - \sqrt{n} \big)$ diverges.} \par\vspace{2cm}
	\end{parts}



% Question 2
\newpage
\question[8] Consider the series given below. 
	\[
	\sum_{n=1}^\infty \dfrac{(-1)}{\sqrt{n}}^{n+1}= 1 - \dfrac{1}{\sqrt{2}} + \dfrac{1}{\sqrt{3}} - \cdots
	\]
	\begin{parts}
	% (a)
	\part Use an appropriate series test to show the series converges. Be sure to fully justify any necessary criterion. \par\vspace{0.3cm}
	\psol{\itshape This is an alternating series. So, we use the Alternating Series Test. Observe\dots} \par
	\psol{\hspace{0.7cm} \itshape \textbullet\ The sequence $\left\{ \tfrac{1}{\sqrt{n}} \right\}$ is decreasing: We need to show that $a_{n+1} \leq a_n$. Observe that $n + 1 \geq n$, so that $\sqrt{n + 1} \geq \sqrt{n}$. But then $\tfrac{1}{\sqrt{n+1}} \leq \tfrac{1}{\sqrt{n}}$. But this shows that $a_{n+1} \leq a_n$. Alternatively, we know that if $f(n)= \tfrac{1}{\sqrt{n}}$, then $f'(n)= -\tfrac{1}{2 n^{3/2}} < 0$ for all $n \geq 1$. But then $f'(n) < 0$ on $[1, \infty)$, so that $f(n)= \tfrac{1}{\sqrt{n}}$ is decreasing.} \par
	\psol{\hspace{0.7cm} \itshape \textbullet\ The limit of the sequence is $0$: $\ds\lim_{n \to \infty} \tfrac{1}{\sqrt{n}}= 0$} \par
	\psol{\itshape Therefore, the series $\ds\sum_{n=1}^\infty \dfrac{(-1)}{\sqrt{n}}^{n+1}$ converges by the Alternating Series Test.} \vfill
	
	% (b)
	\part The sum of the first 48 terms, $S_{48}$, is approximately $0.533106$. What is a maximum possible error bound in this estimation for the sum of the given series? Justify your error bound. \par\vspace{0.3cm}
	\psol{\itshape By the Alternating Series Remainder Theorem, we know the error of the $n$th sum, $S_n$, is at most the $(n+1)$st term, i.e. $a_{n+1}$. So, $|S - S_n|= |\text{error}| \leq a_{n+1}$. Using 48 terms, we know\dots}
		\[
		\psol{|\text{error}| \leq \dfrac{1}{\sqrt{48 + 1}}= \dfrac{1}{\sqrt{49}}= \dfrac{1}{7} \approx 0.142857}
		\] \vfill
	
	\part What is a minimal number of terms one must add to approximate the value of the series to three decimal values, i.e. with an error less than $0.001= \frac{1}{1000}$? Justify your reasoning. \par\vspace{0.3cm}
	\psol{\itshape By the Alternating Series Remainder Theorem, we know the error of the $n$th sum, $S_n$, is at most the $(n+1)$st term, i.e. $a_{n+1}$. So, $|S - S_n|= |\text{error}| \leq a_{n+1}$. So, we can use an $n$ such that $a_{n+1} < \tfrac{1}{1000}$, i.e. $\tfrac{1}{\sqrt{n+1}} < \tfrac{1}{1000}$. But then\dots}
		\[
		\begin{gathered}
		\psol{\tfrac{1}{\sqrt{n + 1}} < \tfrac{1}{1000}} \\
		\psol{\sqrt{n + 1} > 1,\!000} \\
		\psol{n + 1 > 1,\!000,\!000} \\
		\psol{n > 999,\!999}
		\end{gathered}
		\]
	\psol{\itshape Therefore, one need only use at least 1,000,000 terms.}
	\end{parts}



% Question 3
\newpage
\question[8] Determine whether the following series converges or diverges. If the series converges, find the \textit{exact} value of the sum---if possible. \par\vspace{0.3cm} 
	\[
	\sum_{n=1}^\infty \dfrac{(-2)}{3^n}^{n-1}
	\] \par\vspace{0.3cm}

\vsol{\itshape Using the Ratio Test, one has\dots
	\[
	\lim_{n \to \infty} \left| \dfrac{\;\;\dfrac{(-2)^{n}}{3^{n+1}}\;\;}{\dfrac{(-2)^{n-1}}{3^n}} \right|= \lim_{n \to \infty} \dfrac{2^n}{3^{n+1}} \cdot \dfrac{3^n}{2^{n-1}}= \lim_{n \to \infty} \dfrac{2^n}{2^{n-1}} \cdot \dfrac{3^n}{3^{n+1}}= \lim_{n \to \infty} \dfrac{2^n}{2^n\ 2^{-1}} \cdot \dfrac{3^n}{3^n\ 3}= \dfrac{2}{3} < 1
	\]
Alternatively, using the Root Test, one has\dots
	\[
	\lim_{n \to \infty} \left| \dfrac{(-2)^{n-1}}{3^n} \right|^{1/n}= \lim_{n \to \infty} \dfrac{2^{1-1/n}}{3}= \dfrac{2^{1-0}}{3}= \dfrac{2}{3} < 1
	\]
Therefore, whether by the Ratio or Root Test, the given series converges absolutely. Alternatively, observe that the series is alternating. We know the sequence $\left\{ \tfrac{2^{n-1}}{3^n} \right\}$ is decreasing, and $\ds\lim_{n \to \infty} \tfrac{2^{n-1}}{3^n}= 0$. Therefore, the given series converges by the Alternating Series Test. However, neither of these tests finds the sum. \par\vspace{0.3cm}

However, observe\dots
	\[
	\sum_{n=1}^\infty \dfrac{(-2)^{n-1}}{3^n}= \sum_{n=1}^\infty \dfrac{(-2)^n (-2)^{-1}}{3^n}= \sum_{n=1}^\infty -\dfrac{1}{2} \cdot \left( -\dfrac{2}{3} \right)^n
	\]
This series is geometric with $r= -\tfrac{2}{3}$. Therefore, the series $\ds\sum_{n=1}^\infty \dfrac{(-2)^{n-1}}{3^n}$ converges by the Geometric Series Test because $|r|= \tfrac{3}{5} < 1$. Furthermore, we know that\dots
	\[
	\sum_{n=1}^\infty \dfrac{(-2)^{n-1}}{3^n}= \dfrac{\text{first term}}{1 - r}= \dfrac{\tfrac{(-2)^0}{3^1}}{1 - \left(-\tfrac{2}{3} \right)}= \dfrac{\tfrac{1}{3}}{1 + \tfrac{2}{3}} \cdot \dfrac{3}{3}= \dfrac{1}{3 + 2}= \dfrac{1}{5}
	\]
}
	


% Question 4
\newpage
\question[16] For each of the following series, determine whether the series converges or diverges. If the series converges, find the \textit{exact} value of the sum---if possible. \par\vspace{0.3cm}
	\begin{parts}
	% (a)
	\part $\ds\sum_{n=0}^\infty \dfrac{n^{9} \,3^n}{n!}$ \par\vspace{0.3cm}
	\psol{\itshape Using the Ratio Test, we have\dots}
		\[
		\begin{aligned}
		\psol{\lim_{n \to \infty} \left| \dfrac{\;\;\dfrac{(n + 1)^{9} \,3^{n+1}}{(n+1)!}\;\;}{\dfrac{n^{9} \,3^n}{n!}} \right|}&\psol{= \lim_{n \to \infty} \dfrac{(n + 1)^{9} \,3^{n+1}}{(n+1)!} \cdot \dfrac{n!}{n^{9} \,3^n}} \\
		&\psol{= \lim_{n \to \infty} \dfrac{(n+1)^{9}}{n^{9}} \cdot \dfrac{n!}{(n+1)!} \cdot \dfrac{3^{n+1}}{3^n}} \\
		&\psol{= \lim_{n \to \infty} \left( \dfrac{n+1}{n} \right)^{9} \cdot \dfrac{n!}{(n+1)n!} \cdot \dfrac{3^n \,3}{3^n}} \\
		&\psol{= \lim_{n \to \infty} \left( \dfrac{n+1}{n} \right)^{9} \cdot \dfrac{1}{n + 1} \cdot 3} \\
		&\psol{= 1^{9} \cdot 0 \cdot 3} \\
		&\psol{= 0 < 1}
		\end{aligned}
		\]
	\psol{\itshape Therefore, the series $\ds\sum_{n=0}^\infty \dfrac{n^{9} \,3^n}{n!}$ converges absolutely by the Ratio Test. {\tiny. [Note. With techniques we will eventually see, one can show that $\ds\sum_{n=0}^\infty \dfrac{n^{9} \,3^n}{n!}= 5597643 e^3$]}}
	
	% (b)
	\part $\ds\sum_{n=2}^\infty \left( \dfrac{1 - n^2}{5n^2 + 3} \right)^n$ \par\vspace{0.3cm}
	\psol{\itshape Using the Root Test, we have\dots}
		\[
		\psol{\lim_{n \to \infty} \left| \left( \dfrac{1 - n^2}{5n^2 + 3} \right)^n \right|^{1/n}= \lim_{n \to \infty} \left| \dfrac{1 - n^2}{5n^2 + 3}\right|= \left| \dfrac{-1}{5} \right|= \dfrac{1}{5} < 1}
		\]
	\psol{\itshape Therefore, the series $\ds\sum_{n=2}^\infty \left( \dfrac{1 - n^2}{5n^2 + 3} \right)^n$ converges absolutely by the Root Test. [Note. One can show this with the Ratio Test, but it is far more tedious.]}
	\end{parts}



% Question 5
\newpage
\question[16] For each of the following series, determine whether the series converges or diverges. If the series converges, find the \textit{exact} value of the sum---if possible. \par\vspace{0.3cm}
	\begin{parts}
	% (a)
	\part $\ds\sum_{n=7}^\infty \dfrac{\sqrt{n}}{n - 6}$ \par\vspace{0.3cm}
	\psol{\itshape We have\dots}
		\[
		\psol{\lim_{n \to \infty} \dfrac{\;\;\dfrac{1}{\sqrt{n}}\;\;}{\dfrac{\sqrt{n}}{n - 6}}= \lim_{n \to \infty} \dfrac{n - 6}{n}= 1 < \infty \text{ and nonzero}}
		\]
	\psol{\itshape Therefore, the series $\ds\sum_{n=7}^\infty \dfrac{\sqrt{n}}{n - 9}$ and $\ds\sum_{n=7}^\infty \dfrac{1}{\sqrt{n}}$ both converge or both diverge by the Limit Comparison Test. But the series $\ds\sum_{n=7}^\infty \dfrac{1}{\sqrt{n}}$ diverges by the $p$-test with $p= \tfrac{1}{2} \leq 1$. Therefore, the series $\ds\sum_{n=7}^\infty \dfrac{\sqrt{n}}{n - 9}$ diverges. Alternatively, we have\dots}
		\[
		\psol{\sum_{n=7}^\infty \dfrac{\sqrt{n}}{n - 6} \geq \sum_{n=7}^\infty \dfrac{\sqrt{n}}{n}= \sum_{n=7}^\infty \dfrac{1}{\sqrt{n}}}
		\]
	\psol{\itshape The series $\ds\sum_{n=7}^\infty \dfrac{1}{\sqrt{n}}$ diverges by the $p$-test with $p= \tfrac{1}{2} \leq 1$. Therefore, the series $\ds\sum_{n=7}^\infty \dfrac{\sqrt{n}}{n - 6}$ diverges by the Direct Comparison Test. [Note. One could also show the series diverges using the Integral Test. However, we will not give this solution.]} \par\vspace{0.3cm}\vfill
	
	% (b)
	\part $\ds\sum_{n=1}^\infty \sin \left( \dfrac{1}{n^8} \right)$ \par\vspace{0.3cm}
	\psol{\itshape Observe\dots}
		\[
		\psol{\lim_{n \to \infty} \dfrac{\;\;\sin \left( \tfrac{1}{n^8} \right)\;\;}{\dfrac{1}{n^8}}= 1 < \infty \text{ and nonzero}}
		\]
	\psol{\itshape Therefore, the series $\ds\sum_{n=1}^\infty \sin \left( \frac{1}{n^8} \right)$ and $\ds\sum_{n=1}^\infty \frac{1}{n^8}$ both converge or both diverge. The series $\ds\sum_{n=1}^\infty \frac{1}{n^8}$ converges by the $p$-test with $p= 8 > 1$. Therefore, the series $\ds\sum_{n=1}^\infty \sin \left( \dfrac{1}{n^8} \right)$ converges.} \par\vspace{2cm} \vfill
	\end{parts}



% Question 6
\newpage
\question[18] Determine whether the following series converges or diverges. If the series converges, determine whether the series converges conditionally or converges absolutely. \par\vspace{0.3cm}
	\[
	\sum_{n=3}^\infty (-1)^n \; \dfrac{n^{11} - 1}{n^{12} + 5}
	\]

\vsol{\itshape The given series is alternating. Observe\dots \par
	\hspace{0.7cm} \itshape \textbullet\ The sequence $\left\{ \frac{n^{11} - 1}{n^{12} + 5} \right\}$ is decreasing: We need to show that $a_{n+1} \leq a_n$. We shall not show this directly, as the algebra and inequalities are tedious. However, observe that if $f(n)= \frac{n^{11} - 1}{n^{12} + 5}$, then $f'(n)= \tfrac{n^{10}(-n^{12} + 12n + 55)}{(n^{12} + 5)^2}$, which is negative for $n > 1$, so that $f(n)$ is decreasing. \par
	\hspace{0.7cm} \itshape \textbullet\ The limit of the sequence is $0$: $\ds\lim_{n \to \infty} \frac{n^{11} - 1}{n^{12} + 5}= 0$.

Therefore, the series $\ds\sum_{n=3}^\infty (-1)^n \; \dfrac{n^{11} - 1}{n^{12} + 5}$ converges by the Alternating Series Test. 

Now we need to determine whether the series $\ds\sum_{n=3}^\infty \dfrac{n^{11} - 1}{n^{12} + 5}$ converges. Observe\dots
	\[
	\lim_{n \to \infty} \dfrac{\;\;\dfrac{n^{11} - 1}{n^{12} + 5}\;\;}{\dfrac{1}{n}}= \lim_{n \to \infty} \dfrac{n^{11} - 1}{n^{12} + 5} \cdot \dfrac{n}{1}= \lim_{n \to \infty} \dfrac{n^{12} - n}{n^{12} + 5}= 1 < \infty \text{ and nonzero}
	\]
Therefore, the series $\ds\sum_{n=3}^\infty \dfrac{n^{11} - 1}{n^{12} + 5}$ and $\ds\sum_{n=3}^\infty \dfrac{1}{n}$ both converge or both diverge by the Limit Comparison Test. The series $\ds\sum_{n=3}^\infty \dfrac{1}{n}$ diverges by the $p$-test with $p= 1 \leq 1$ (we also know that the Harmonic Series diverges). Therefore, the series $\ds\sum_{n=3}^\infty \dfrac{n^{11} - 1}{n^{12} + 5}$ diverges. Alternatively, because $1 \leq \frac{n}{2} < \frac{n^{11}}{2}$ for $n \geq 2$, we have\dots
	\[
	\sum_{n=3}^\infty \dfrac{n^{11} - 1}{n^{12} + 5} \geq \sum_{n=3}^\infty \dfrac{n^{11} - \tfrac{n^{11}}{2}}{n^{12} + 5n^{12}}= \sum_{n=3}^\infty \dfrac{\tfrac{n^{11}}{2}}{6n^{14}}= \dfrac{1}{12} \sum_{n=3}^\infty \dfrac{1}{n}
	\]
so that $\ds\sum_{n=3}^\infty \dfrac{n^{11} - 1}{n^{12} + 5}$ diverges by the Direct Comparison Test. 

Therefore, the series $\ds\sum_{n=3}^\infty (-1)^n \; \dfrac{n^{11} - 1}{n^{12} + 5}$ is conditionally convergent.}



% Question 7
\newpage
\question[18] Use the Integral Test to determine whether the following series converges or diverges. If the series converges, find the \textit{exact} value of the sum---if possible. \par\vspace{0.3cm}
	\[
	\sum_{n=2}^\infty \dfrac{1}{n \ln n}
	\]

\vsol{\itshape Let $f(x)= \tfrac{1}{x \ln x}$. We know\dots \par\vspace{0.3cm}
	\hspace{0.7cm} \itshape \textbullet\ $f(x)$ is positive: $x > 0$ and $\ln x > 0$ for $x > 1$, so that $x \ln x > 0$, which shows that $f(x)= \tfrac{1}{x \ln x} > 0$. \par
	\hspace{0.7cm} \itshape \textbullet\ $f(x)$ is decreasing: We know that $f'(x)= - \dfrac{\ln x + 1}{(x \ln x)^2} < 0$ for $x > \frac{1}{e}$. \par
	\hspace{0.7cm} \itshape \textbullet\ $f(x)$ is continuous: We know that $1$, $x$, and $\ln x$ are continuous for $x > 0$. Therefore, $x \ln x$ is continuous for $x > 0$---being a product of continuous functions. We know $x \ln x= 0$ if and only if $x= 0$ or $x= 1$. But then $f(x)= \frac{1}{x \ln x}$ is continuous for $x > 1$---being a quotient of continuous functions. \par\vspace{0.3cm}

Therefore, the Integral Test applies. So, $\ds\sum_{n=2}^\infty \dfrac{1}{n \ln n}$ converges if and only if $\ds\int_2^\infty \dfrac{1}{x \ln x} \;dx$ converges. \par\vspace{0.3cm}

Let $u= \ln x$, so that $du= \frac{1}{x}\;dx$. But then\dots
	\[
	\int \dfrac{1}{x \ln x} \;dx= \int \dfrac{1}{\ln x} \cdot \dfrac{1}{x} \;dx= \int \dfrac{1}{u} \;du= \ln|u| + C= \ln|\ln x| + C
	\]
But then\dots
	\[
	\int_2^\infty \dfrac{1}{x \ln x} \;dx= \lim_{b \to \infty} \int_2^b \dfrac{1}{x \ln x} \;dx= \lim_{b \to \infty}  \ln|\ln x| \bigg|_2^b= \lim_{b \to \infty} \ln| \ln b | - \ln(\ln(2))= \infty
	\] \par\vspace{0.3cm}
Therefore, $\ds\sum_{n=2}^\infty \dfrac{1}{n \ln n}$ diverges.}

\end{questions}
\end{document}